\chapter{绪论}

\section{研究背景}

温度历程决定金属热处理中的相变驱动力、组织均匀性和热应力水平。炉温可以直接测量和调节，工件内部温度场通常只能通过少量热电偶间接获得。厚壁工件升温时，表面与心部会形成明显温差。控制系统若只依据炉温或表面温度，容易错判心部到温时刻。因此，能够随控制输入快速推演完整温度场的模型具有直接工程价值。

有限差分、有限体积和有限元方法已经构成热过程分析的标准工具。它们能够处理温变物性和复杂边界，也适合生成高可信度参考解。在线控制提出了另一类计算需求。控制器需要在每个决策时刻比较多条候选炉温序列，并根据新观测重新计算。二维、三维模型中的重复求解成本很高。学习型代理模型可将主要计算转移到离线阶段，为实时推演提供条件。

现有温度场机器学习研究多采用两种任务设置。第一类直接近似给定偏微分方程的解 $T(x,t)$。物理信息神经网络通过自动微分把控制方程和边界条件写入训练目标，在正问题和逆问题中形成了统一框架\cite{raissi2019pinn,karniadakis2021piml}。第二类从有限测点恢复当前或历史温度场。相关方法已用于钢材加热、水淬全场重构和虚拟热传感器\cite{sun2025steel,zhao2022quenching,go2023sensor}。这些任务推动了温度场学习，但工艺决策还需要一个动作条件化的状态转移模型。模型应接收当前温度场和下一段炉温指令，并能递归计算控制改变后的未来状态。

世界模型提供了合适的动力学表述。本文采用工程化定义，将世界模型写为
\begin{equation}
  \bm{s}_{n+1}=F_{\theta}(\bm{s}_n,a_n,\bm{p}),
  \label{eq:wm_definition}
\end{equation}
其中 $\bm{s}_n$ 是离散温度场，$a_n$ 是炉温动作，$\bm{p}$ 包含几何、物性和边界参数。式\eqref{eq:wm_definition}可以连续调用，所以同一模型能够评价多条候选工艺轨迹。该定义保留了热传导状态的物理含义，也便于与模型预测控制连接。

\section{相关研究}

\subsection{物理信息学习与温度场预测}

PINN 将方程残差、初边值条件和观测误差共同用于网络训练\cite{raissi2019pinn}。它减少了对标注数据的依赖，并可在同一计算图中处理未知参数。训练效果会受到损失尺度、梯度冲突和采样分布影响\cite{wang2021gradient,krishnapriyan2021failure}。因此，物理残差较低并不自动对应更小的场误差。温度场应用还要面对长时递归误差、控制分布变化和边界参数漂移。评价口径需要覆盖完整轨迹。

Sun 等将物理信息网络用于热轧钢加热温度分布预测\cite{sun2025steel}。Zhao 等利用稀疏测温信息恢复水淬过程的时空温度场，并估计热边界\cite{zhao2022quenching}。Go 等研究了有限温度观测下的虚拟热传感器和未知热流估计\cite{go2023sensor}。这些工作证明了物理信息学习在金属热过程中的可行性。本文关注控制输入驱动的未来轨迹。训练和评价均采用闭环递归方式，模型输出直接进入后续状态和滚动优化。

\subsection{神经算子与学习型动力学模型}

DeepONet 和 Fourier Neural Operator 学习输入函数或参数到解场之间的映射\cite{lu2021deeponet,li2021fno}。这类模型适合跨网格、跨载荷或跨几何问题。本文当前对象为固定的一维 41 节点网格。轻量残差多层感知机已经能够覆盖状态维数，并具有较低推理开销。研究重点因此放在动力学训练、物理约束和边界可辨识性上。

学习型物理模拟器常采用自回归推演。一步监督训练只接触真实状态，部署时却持续接收模型生成状态，两种输入分布的偏移会放大长期误差。List 等的研究表明，时间展开方式会显著影响瞬态动力学的轨迹精度\cite{list2025unrolled}。本文使用五步闭环展开，并以完整验证轨迹选择检查点。一步误差只作为诊断指标。

\subsection{动态热边界辨识}

高温炉内钢材加热同时受对流和辐射影响。辐射贡献随表面温度快速变化，钢材氧化状态还会改变发射率。Kim 的加热炉模型说明了辐射边界和钢坯内部瞬态导热的耦合关系\cite{kim2007furnace}。边界参数往往无法直接测量，逆热传导方法需要从表面温度或热流响应中估计载荷。Frankel 和 Arimilli 研究了瞬态表面测温下对流与辐射载荷的反演\cite{frankel2007loads}。序贯正则化和神经正向模型也已用于在线热流估计\cite{chen2019online,mohebbi2021htc}。

对流系数 $h(t)$ 和发射率 $\varepsilon(t)$ 若同时按时间自由变化，单个时间步的温度响应无法提供两个独立观测方向。将二者直接作为网络输入，会把等价热边界表示成多组不同参数。本文对离散边界方程进行灵敏度分析，确定可观测组合 $\Heff(t)$，并据此设计世界模型输入和在线状态估计器。

\section{研究内容与贡献}

本文围绕受控热处理过程建立一条完整方法链。主要贡献如下。

\begin{enumerate}
  \item 提出面向金属热处理温度场的受控递归世界模型。模型显式接收炉温动作和热参数，采用五步闭环展开训练。完整轨迹 RMSE、尾部误差和物理一致性共同构成评价标准。
  \item 建立覆盖控制曲线、换热系数、发射率和温变物性的 OOD 实验矩阵。结果给出物理约束的适用条件。其主要收益集中在组合参数外推、冷却轨迹、尾部误差和离散能量一致性。
  \item 推导动态对流辐射边界的可辨识表示 $\Heff$。基于该表示构建因果观测器、集合卡尔曼滤波和滚动控制。独立 BDF 求解器、稀疏测点和风险感知控制实验验证了方法链的有效性。
\end{enumerate}

全文共六章。第二章给出物理模型和数值基准。第三章介绍世界模型、物理损失和边界参数化。第四章分析多步推演、参数 OOD 与跨求解器结果。第五章研究动态边界辨识、部分可观测状态估计和闭环控制。第六章总结研究结论并给出后续研究方向。
